Put $b_k=\|S^{(k)}f_k\|_H$, $h=(\sum_k b_k^2)^{1/2}$, and $W=(\sum_k |f_k|^2)^{1/2}$. Let $M$ be the Hardy--Littlewood maximal operator over cubes and set $L=\lceil(N+2)\log_2(2/\delta)\rceil$. For each dyadic cube $Q$ of side $2^{-j}$ and center $y$, compare $h$ with $c_Q=(\sum_{k\le j-L} b_k(y)^2)^{1/2}$. For $x,z\in Q$ and $k\le j-L$, the derivative bound and the mean value inequality give
\[
 \|S^{(k)}f_k(x)-S^{(k)}f_k(y)\|_H
 \lesssim C_*\delta^{-N}2^{k-j}MW(z).
\]
Indeed $2^k|x-z|\lesssim1$, and shifted kernels on the segment from $y$ to $x$ are dominated by $2^{nk}(1+2^k|z-\cdot|)^{-N}$ up to a fixed constant. Convolving this radial decreasing kernel with $W$ costs at most a constant times $MW(z)$. The square sum of the differences costs $CC_*\delta MW(z)$ since $\delta^{-N}2^{-L}\lesssim\delta$.

For $k\ge j+L$ split $f_k$ on $5Q$ and its complement, where $5Q$ is the concentric dilation of $Q$. Dilation preserves the $L^2$ operator norm, so for $z\in Q$,
\[
 \frac1{|Q|}\int_Q\sum_{k\ge j+L}
 \|S^{(k)}(f_k\mathbf1_{5Q})(x)\|_H^2\,dx
 \lesssim B_2^2 M(W^2)(z).
\]
For $x,z\in Q$, the kernel bound summed on successive doubles of $5Q$ gives
\[
 \|S^{(k)}(f_k\mathbf1_{(5Q)^c})(x)\|_H
 \lesssim C_*\delta^{-N}2^{-(N-n)(k-j)} MW(z).
\]
On a cube of side comparable to $2^{a-j}$, the distance from $x$ to the associated exterior shell is bounded below by a constant times $2^{a-j}$. The integral of $|f_k|$ over the cube is bounded by its volume times $MW(z)$, leaving a convergent sum in $a$ because $N>n$. The square sum over $k\ge j+L$ then costs $CC_*\delta MW(z)$.

For the fewer than $2L$ remaining $k$, H\"older in this finite set gives
\[
 \left(|Q|^{-1}\int_Q\sum_{|k-j|<L}b_k^2\right)^{1/2}
 \le(2L)^{1/2-1/p}F(z),\qquad
 F=\left(\sum_k[M(b_k^2)]^{p/2}\right)^{1/p}.
\]
As $p/2>1$, the scalar maximal theorem and invariance of $B_p$ under dilation yield
\[
 \|F\|_p^p\lesssim\sum_k\|b_k\|_p^p
 \le B_p^p\sum_k\|f_k\|_p^p\le B_p^p\|W\|_p^p.
\]
The triangle inequality for the Hilbert square norm and Cauchy--Schwarz on $Q$ now bound the mean of $|h-c_Q|$ by
\[
 C\left(C_*\delta MW(z)+B_2M(W^2)(z)^{1/2}
       +L^{1/2-1/p}F(z)\right).
\]
Replacing $c_Q$ by the mean $h_Q$ costs at most two. The dyadic sharp function $h^\#(z)=\sup_{Q\ni z}|Q|^{-1}\int_Q|h-h_Q|$, with the supremum over dyadic cubes, therefore has its $L^p$ norm bounded by the stated right-hand side, using the maximal theorem for $p$ and $p/2$. Apply the sharp function inequality $\|h\|_p\lesssim_p\|h^\#\|_p$ from \cite{r-74}, Chapter~IV, Section~3.6, Corollary~1. Its smaller-exponent hypothesis holds because
\[
 \|h\|_2^2=\sum_k\|S^{(k)}f_k\|_{L^2(H)}^2
 \le B_2^2\sum_k\|f_k\|_2^2<\infty
\]
and $2<p$.
